\chapter{Comparison with Fréchet calculus}
\label{chap:frechet}

On normed spaces two calculi are available: the classical Fréchet calculus, in which the
derivative is required to depend continuously on the point in the operator norm, and the
calculus of Michal and Bastiani, which only asks for joint continuity of
$(x, v) \mapsto \dd f(x)(v)$. The second notion is weaker; Fréchet $\Ck{k}$ maps are
Michal--Bastiani $\Ck{k}$ (Proposition~\ref{prop:frechet-to-mb}). In this chapter we prove the
converse, up to the loss of one derivative. As a consequence both notions of smoothness agree,
so that for smooth maps between normed spaces one may use whichever calculus is more
convenient.

In the second part of the chapter we apply this comparison to maps with bounded below
derivative, and show that such maps reflect smoothness. Throughout, $E$ and $F$ are real normed
spaces; completeness is assumed only where it is stated.

\section{From Michal--Bastiani to Fréchet}

\begin{theorem}[Michal--Bastiani $\Ck{k+1}$ implies Fréchet $\Ck{k}$]
  \label{thm:mb-to-frechet}
  \lean{LocallyConvex.Analysis.ContDiffMBOn.contDiffOn, LocallyConvex.Analysis.ContDiffMBOn.contDiffOn_one}
  \leanok
  \uses{thm:linear, lem:mean-value-inclusion, lem:continuity}
  If $f$ is $\Ck{k+1}$ on $U$ in the sense of Michal--Bastiani, then $f$ is Fréchet $\Ck{k}$ on
  $U$.
\end{theorem}

The extra derivative is what converts joint continuity into continuity in the operator norm:
in the proof below, joint continuity of the next derivative, together with the mean value
inclusion, shows that the derivatives are locally Lipschitz in the operator norm.

\begin{proof}
  \leanok
  We prove, by induction on $k$ and for all $m$ simultaneously, the following statement:
  \emph{if $g \colon U \to \Lin^m(E; F)$ is such that
  $\hat g(x, v_1, \dots, v_m) = g(x)(v_1, \dots, v_m)$ is $\Ck{k+1}$ in the sense of
  Michal--Bastiani, then $g$ is Fréchet $\Ck{k}$}. Here $\Lin^m(E; F)$ is the space of continuous
  $m$-linear maps with the operator norm. For $m = 0$ this is the theorem.

  The candidate for the derivative of $g$ at $y$ is the $(m+1)$-linear map
  $B_y(w, v) = \dd \hat g\bigl((y, v), (w, 0)\bigr)$, the derivative of $\hat g$ in a base
  direction; it is multilinear by Theorem~\ref{thm:linear}. Since $\dd \hat g$ is jointly
  continuous and vanishes in the direction $0$, the maps $B_{y'}$ are uniformly bounded in
  operator norm for $y'$ near $y$. With the mean value inclusion
  (Lemma~\ref{lem:mean-value-inclusion}), this shows that $g$ is locally Lipschitz into
  $\Lin^m(E; F)$; in particular it is continuous, which is the case $k = 0$.

  For the inductive step, suppose that $\hat g$ is $\Ck{k+2}$. Then
  $(y, w, v) \mapsto B_y(w, v)$ is $\Ck{k+1}$, so by the case $(k, m + 1)$ the
  map $y \mapsto B_y$ is Fréchet $\Ck{k}$, and in particular continuous in operator norm. The mean
  value inclusion then gives
  \[
    \|g(y + w) - g(y) - B_y(w, \cdot)\| = o(\|w\|),
  \]
  so $g$ is Fréchet differentiable with derivative $y \mapsto B_y$, and $g$ is Fréchet
  $\Ck{k+1}$.
\end{proof}

\begin{remark}
  This is \cite[2.3.17(b)]{GN26}. The proof in \cite{GN26} goes through the compact-open topology on
  $\Lin(E, F)$ and the exponential law (\cite[1.7.30]{GN26}). The proof above stays in normed spaces:
  the derivatives are tracked as continuous multilinear maps with the operator norm, which is
  closer to the existing Fréchet calculus of Mathlib. The price is the auxiliary statement for
  all $m$, which has to be proved simultaneously.
\end{remark}

Since the loss of one derivative disappears when $k = \infty$, we obtain the following.

\begin{corollary}[Smooth maps agree]
  \label{cor:smooth-iff}
  \lean{LocallyConvex.Analysis.smoothMBOn_iff_contDiffOn}
  \leanok
  \uses{thm:mb-to-frechet, prop:frechet-to-mb}
  A map between open subsets of normed spaces is smooth in the sense of Michal--Bastiani if and
  only if it is smooth in the sense of Fréchet.
\end{corollary}

\begin{proof}
  \leanok
  Theorem~\ref{thm:mb-to-frechet} and Proposition~\ref{prop:frechet-to-mb}.
\end{proof}

\section{Operator-valued maps on locally convex spaces}

\begin{theorem}[Operator-valued maps]
  \label{thm:operator-valued}
  \lean{LocallyConvex.Analysis.hasDirDerivAt_slice_of_unc_clm,
    LocallyConvex.Analysis.dMB_unc_clm_smul, LocallyConvex.Analysis.dMB_unc_clm_add,
    LocallyConvex.Analysis.hasDerivAt_line_of_unc_clm,
    LocallyConvex.Analysis.exists_bound_unc_clm,
    LocallyConvex.Analysis.continuousOn_of_unc_clm,
    LocallyConvex.Analysis.hasDirDerivAt_of_unc_clm,
    LocallyConvex.Analysis.contDiffMBOn_of_unc_clm}
  \leanok
  \uses{thm:mb-to-frechet}
  Let $E$ be a locally convex space, $F_1$, $F_2$ normed spaces, $U \subseteq E$ open and
  $g \colon U \to \Lin(F_1, F_2)$. If $\hat g(y, \xi) = g(y)\xi$ is $\Ck{k+1}$ on $U \times F_1$
  in the sense of Michal--Bastiani, then $g$ is $\Ck{k}$ for the operator norm on
  $\Lin(F_1, F_2)$.
\end{theorem}

\begin{proof}
  \leanok
  Let $D(y, \xi, w) = \dd\hat g((y, \xi), (w, 0))$; it is linear in $\xi$ and in $w$, and
  continuous. Continuity at $(y_0, 0, 0)$ gives a convex neighborhood $N$ of $y_0$, a
  neighborhood $W$ of $0$ and $C$ with $\|D(y, \xi, w)\| \le C\|\xi\|$ for $y \in N$,
  $w \in W$. If $y - y_0 \in rW$, the mean value inequality along the segment
  $[y_0, y] \subseteq N$ gives $\|g(y) - g(y_0)\| \le rC$: $g$ is continuous. The operator
  $h(y, w) = D(y, \cdot, w)$ is the directional derivative of $g$ for the operator norm, by the
  mean value inequality applied to $s \mapsto g(y + sw)\xi - s\,h(y, w)\xi$ and the continuity of
  $h$, which is the case $k = 0$ applied to $h$ on $U \times E$. The induction runs on the domain
  $E \times E$ of the derivative.
\end{proof}

The statement is meaningless if $F_1$ is not normable: no vector topology on $\Lin(F_1, F_2)$
makes the evaluation continuous. Theorem~\ref{thm:mb-to-frechet} is the case where $E$ is normed;
there the multilinear version is needed for the induction.

\section{Maps with bounded below derivative}

In this section we consider a Banach space $E$ (complete), a normed space $F$ (complete or not),
and a map $j \colon U \to F$, smooth on an open set $U \subseteq E$, with derivative bounded
below on $U$: there is $m > 0$ with
\[
  m \|v\| \le \|\dd j(x)(v)\| \qquad (x \in U,\ v \in E).
\]
By Corollary~\ref{cor:smooth-iff} it makes no difference in which sense $j$ is smooth. No
complement of the range of $\dd j(x)$ is assumed. The simplest example is a continuous linear map
$j \colon E \to F$ with $m\|v\| \leq \|j(v)\|$, for instance an isometric embedding; in that
case $\dd j(x) = j$ for every $x$.

Our aim is to show that such a $j$ reflects smoothness: a continuous map $\gamma$ into $U$ is
of class $\Ck{k}$ as soon as $j \circ \gamma$ is. The test maps $\gamma$ are defined on open
subsets of arbitrary locally convex spaces. We first prove a linear version, in which the
unknown appears linearly.

\begin{lemma}[Linear bootstrap]
  \label{lem:bootstrap-linear}
  \lean{LocallyConvex.Analysis.ContDiffMBOn.of_fderiv_eq_of_boundedBelow}
  \leanok
  \uses{thm:mb-to-frechet, cor:smooth-iff}
  Let $W$ be an open subset of a locally convex space $H$, $\gamma \colon W \to U$ and
  $w \colon W \to F$ of class $\Ck{k}$, and $z \colon W \to E$ continuous with
  $\dd j(\gamma q)(z q) = w(q)$ for $q \in W$. Then $z$ is of class $\Ck{k}$.
\end{lemma}

\begin{proof}
  \leanok
  By induction on $k$. For the inductive step, fix $q \in W$ and a direction $h$, and write
  $A_t = \dd j(\gamma(q + th))$. From the equation,
  \[
    A_0\Bigl(\frac{z(q + th) - z(q)}{t}\Bigr)
      = \frac{w(q + th) - w(q)}{t} - \frac{A_t - A_0}{t}\, z(q + th).
  \]
  Since $j$ is Fréchet $C^2$ (Theorem~\ref{thm:mb-to-frechet}), $t \mapsto A_t$ is
  differentiable in the operator norm, so the right-hand side converges. As $A_0$ is bounded
  below and $E$ is complete, the difference quotients of $z$ are Cauchy, hence converge to some
  $D$ with $A_0 D = \dd w(q)(h) - \dd^2 j(\gamma q)(\dd\gamma(q)(h), z q)$.

  The derivative $\dd z$ is continuous, because solutions of $A(p) x(p) = r(p)$ are continuous
  when $A$ is continuous in the operator norm and uniformly bounded below. Moreover, it
  satisfies an equation of the same type on $W \times H$, with a right-hand side of class
  $\Ck{k}$ by the inductive hypothesis applied to $z$. So $\dd z$ is $\Ck{k}$ by the inductive
  hypothesis again, and $z$ is $\Ck{k+1}$.
\end{proof}

\begin{theorem}[Bounded below derivatives reflect smoothness]
  \label{thm:bounded-below}
  \lean{LocallyConvex.Analysis.ContDiffMBOn.of_comp_of_boundedBelow}
  \leanok
  \uses{lem:bootstrap-linear}
  Let $W$ be an open subset of a locally convex space and $\gamma \colon W \to U$ continuous. If
  $j \circ \gamma$ is of class $\Ck{k}$, then so is $\gamma$.
\end{theorem}

In the linear example above, the theorem says that a continuous map into $E$ is $\Ck{k}$ as soon
as it is $\Ck{k}$ when regarded, through $j$, as a map into the larger space $F$.

\begin{proof}
  \leanok
  By induction on $k$. Fix $q \in W$, a direction $h$, and write $A = \dd j(\gamma q)$,
  $d_t = \gamma(q + th) - \gamma(q)$. Since $j$ is strictly differentiable at $\gamma q$, for
  any $c > 0$ we have
  \[
    \|j(\gamma(q + th)) - j(\gamma q) - A d_t\| \le c \|d_t\|
  \]
  for small $t$. With $c \le m/2$ the lower bound gives
  $\|d_t / t\| \le (2/m) \|(j(\gamma(q+th)) - j(\gamma q))/t\|$, so
  $A(d_t/t) \to \dd(j\circ\gamma)(q)(h)$, and as in Lemma~\ref{lem:bootstrap-linear} the
  quotients $d_t/t$ converge. The derivative $\dd\gamma$ then solves
  \[
    \dd j(\gamma q)\bigl(\dd\gamma(q)(h)\bigr) = \dd(j \circ \gamma)(q)(h),
  \]
  and Lemma~\ref{lem:bootstrap-linear} applies.
\end{proof}

\begin{remark}
  A word of caution on the scope of this result. The proof uses that $x \mapsto \dd j(x)$ is
  continuous in the operator norm, to keep the lower bound near a point. The Michal--Bastiani
  definition only gives the joint continuity of $(x, v) \mapsto \dd j(x)(v)$; on normed spaces
  the operator norm continuity is recovered through Theorem~\ref{thm:mb-to-frechet}. For
  locally convex model spaces, Section~\ref{sec:equi-embedding} proves the statement under a
  locally uniform lower bound, which is automatic in normed spaces; with only a pointwise lower
  bound it is open.
\end{remark}

\section{Uniformly injective derivatives on locally convex spaces}
\label{sec:equi-embedding}

We now let $E$, $F$ be Hausdorff locally convex spaces, $E$ sequentially complete
(Definition~\ref{def:seq-complete}), and $j \colon U \to F$ of class $\Ck{2}$ on an open set
$U \subseteq E$. For a single continuous linear map $L \colon E \to F$ the analogue of being
bounded below is being a topological embedding: for every continuous seminorm $p$ on $E$ there
are a continuous seminorm $q$ on $F$ and $C \ge 0$ with $p(w) \le C\, q(Lw)$. The proof of
Theorem~\ref{thm:bounded-below} needs this near each point, uniformly, and also for convex
combinations of derivatives. We found no source for the results of this section; they were
written for this library by adapting the Banach proof, and we do not claim they are new.

\begin{definition}[Uniformly injective derivatives]
  \label{def:equi-embedding}
  \lean{LocallyConvex.Analysis.EquiEmbeddingOn, LocallyConvex.Analysis.EquiEmbeddingNear}
  \leanok
  The derivatives of $j$ are \emph{uniformly injective on} $V \subseteq U$ if for every
  continuous seminorm $p$ on $E$ there are a continuous seminorm $q$ on $F$ and $C \ge 0$ with
  \[
    p(w) \le C\, q(z) \qquad \text{for all } w \in E \text{ and }
    z \in \overline{\operatorname{conv}}\,\{\dd j(y)(w) : y \in V\}.
  \]
  We say that this holds \emph{near} $x_0$ if it holds on some neighborhood of $x_0$.
\end{definition}

Equivalently, $p(w) \le C\, q\bigl(\sum_i \lambda_i \dd j(y_i)(w)\bigr)$ for all finite convex
combinations with $y_i \in V$. Taking a single point, each $\dd j(y)$, $y \in V$, is a
topological embedding with constants independent of $y$. The hypothesis does not follow from
each $\dd j(x)$ being an embedding: in the Fr\'echet space of rapidly decreasing sequences
there is a smooth family of linear automorphisms $L(t)$ with $L(0) = \mathrm{id}$ whose inverses
are not equicontinuous near $t = 0$
(\texttt{LocallyConvex.Analysis.Counterexamples.RapidDecay.counterexample}).

\begin{lemma}[The normed case]
  \label{lem:equi-embedding-normed}
  \lean{LocallyConvex.Analysis.equiEmbeddingNear_of_boundedBelow}
  \leanok
  \uses{def:equi-embedding, thm:mb-to-frechet}
  If $E$, $F$ are normed and $\dd j(x_0)$ is bounded below, then the derivatives of $j$ are
  uniformly injective near $x_0$.
\end{lemma}

\begin{proof}
  \leanok
  Let $m\|v\| \le \|\dd j(x_0) v\|$. By Theorem~\ref{thm:mb-to-frechet}, $x \mapsto \dd j(x)$
  is continuous in the operator norm, so $\|\dd j(y) - \dd j(x_0)\| < m/2$ on a neighborhood
  $V$ of $x_0$. The closed ball of radius $\frac m2\|w\|$ around $\dd j(x_0) w$ is closed and
  convex and contains every $\dd j(y) w$, $y \in V$, hence every $z$ in their closed convex
  hull, and then $\|z\| \ge \frac m2 \|w\|$. A continuous seminorm satisfies $p \le c\|\cdot\|$.
\end{proof}

\begin{lemma}[Difference quotients]
  \label{lem:equi-embedding-quotient}
  \lean{LocallyConvex.Analysis.ContDiffMBOn.exists_tendsto_of_equiEmbedding}
  \leanok
  \uses{def:equi-embedding, lem:mean-value-inclusion, thm:linear}
  Assume the derivatives of $j$ are uniformly injective near $x \in U$. If $d_i \to 0$,
  $\tau_i \ne 0$ and $\tau_i^{-1}(j(x + d_i) - j(x)) \to R$ along a countably generated filter,
  then $\tau_i^{-1} d_i \to D$ with $\dd j(x) D = R$.
\end{lemma}

\begin{proof}
  \leanok
  Write $g_i = \tau_i^{-1} d_i$ and $\Delta_i = \tau_i^{-1}(j(x + d_i) - j(x))$. By the mean
  value inclusion (Lemma~\ref{lem:mean-value-inclusion}), $\Delta_i$ lies in the closed convex
  hull of $\{\dd j(y)(g_i) : y \in [x, x + d_i]\}$, and the segment lies in $V$ for large $i$; so
  $p(g_i) \le C q(\Delta_i)$ is eventually bounded, for each continuous seminorm $p$.

  Joint continuity of $\dd^2 j$ at $(x, 0, 0)$ and homogeneity give, for each continuous
  seminorm $q$ on $F$, a continuous seminorm $r$ on $E$ with
  $q(\dd^2 j(z)(u, v)) \le r(u) r(v)$ for $z$ near $x$. The mean value inclusion, twice, gives
  $q(\Delta_i - \dd j(x) g_i) \le r(d_i)\, r(g_i) \to 0$. Hence $\dd j(x) g_i \to R$. As
  $\dd j(x)$ is an embedding, $(g_{i_n})$ is Cauchy along a sequence $i_n$ tending to the filter,
  and converges to some $D$ because $E$ is sequentially complete; then $\dd j(x) D = R$, and
  $p(g_i - D) \le C q(\dd j(x) g_i - R) \to 0$.
\end{proof}

\begin{lemma}[Linear bootstrap, locally convex case]
  \label{lem:bootstrap-linear-lc}
  \lean{LocallyConvex.Analysis.ContDiffMBOn.of_dMB_eq_of_equiEmbedding}
  \leanok
  \uses{lem:equi-embedding-quotient, lem:glockner}
  Assume the derivatives of $j$ are uniformly injective near every point of $U$ and $j$ is
  $\Ck{k+1}$. Let $W$ be an open subset of a locally convex space, $\gamma \colon W \to U$ and
  $w \colon W \to F$ of class $\Ck{k}$, and $z \colon W \to E$ continuous with
  $\dd j(\gamma q)(z q) = w(q)$. Then $z$ is of class $\Ck{k}$.
\end{lemma}

\begin{proof}
  \leanok
  As in Lemma~\ref{lem:bootstrap-linear}. The second order term
  $t^{-1}\bigl(\dd j(\gamma(q + th))(z(q+th)) - \dd j(\gamma q)(z(q + th))\bigr)$ converges to
  $\dd^2 j(\gamma q)(\dd\gamma(q)(h), z q)$ by Glöckner's lemma (Lemma~\ref{lem:glockner}) for
  $\dd j$, and the limit of the quotients of $z$ is produced as in
  Lemma~\ref{lem:equi-embedding-quotient}. Continuity of $\dd z$: if
  $\dd j(\gamma a)(z a) = r(a)$, then
  $p(z a' - z a) \le C q\bigl(r(a') - \dd j(\gamma a')(z a)\bigr) \to 0$ by joint continuity
  of $\dd j$.
\end{proof}

\begin{theorem}[Uniformly injective derivatives reflect smoothness]
  \label{thm:equi-embedding}
  \lean{LocallyConvex.Analysis.ContDiffMBOn.of_comp_of_equiEmbedding, LocallyConvex.Analysis.ContDiffMBOn.of_comp_of_equiEmbedding_of_smooth}
  \leanok
  \uses{lem:bootstrap-linear-lc, lem:equi-embedding-quotient}
  Assume the derivatives of $j$ are uniformly injective near every point of $U$, and that $j$ is
  $\Ck{2}$ and $\Ck{k}$. Let $W$ be an open subset of a locally convex space and
  $\gamma \colon W \to U$ continuous. If $j \circ \gamma$ is of class $\Ck{k}$, then so is
  $\gamma$.
\end{theorem}

\begin{proof}
  \leanok
  By induction on $k$. Lemma~\ref{lem:equi-embedding-quotient}, with
  $d_t = \gamma(q + th) - \gamma(q)$ and $\tau_t = t$, gives the derivative $\dd\gamma$, which
  solves $\dd j(\gamma q)(\dd\gamma(q)(h)) = \dd(j \circ \gamma)(q)(h)$; then
  Lemma~\ref{lem:bootstrap-linear-lc} applies.
\end{proof}

\begin{corollary}[Pointwise lower bounds in Banach spaces]
  \label{cor:bounded-below-pointwise}
  \lean{LocallyConvex.Analysis.ContDiffMBOn.of_comp_of_boundedBelow_pointwise}
  \leanok
  \uses{thm:equi-embedding, lem:equi-embedding-normed}
  Theorem~\ref{thm:bounded-below} holds when $\dd j(x)$ is bounded below at every $x \in U$, with
  a constant depending on $x$.
\end{corollary}

\begin{proof}
  \leanok
  Lemma~\ref{lem:equi-embedding-normed} and Theorem~\ref{thm:equi-embedding}; a Banach space is
  sequentially complete.
\end{proof}

Whether Theorem~\ref{thm:equi-embedding} holds when each $\dd j(x)$ is only assumed to be a
topological embedding is open.

\section{Quadratic second order equations}

On a Banach space $E$ we consider the equation of the geodesics of a spray,
\[
  c'' = \Gamma(c)(c', c'),
\]
where $\Gamma \colon U \times E \times E \to E$ is $\Ck{2}$ in the sense of Michal--Bastiani on
an open set $U \subseteq E$ and bilinear in the last two variables. It is equivalent to the first
order system $(c, v)' = (v, \Gamma(c)(v, v))$, whose right hand side is Fr\'echet $\Ck{1}$
(Theorem~\ref{thm:mb-to-frechet}), hence locally Lipschitz.

\begin{theorem}[Local existence, uniform for bounded velocities]
  \label{thm:quadratic-ode-existence}
  \lean{LocallyConvex.Analysis.exists_uniform_isSprayCurveOn_of_norm_le}
  \leanok
  \uses{thm:mb-to-frechet}
  Let $E$ be a Banach space and $x_0 \in U$. There is $r > 0$ such that for every $K \ge 0$ there
  is $\varepsilon > 0$ for which every initial condition $c(0) = y$, $c'(0) = w$ with
  $\|y - x_0\| \le r$ and $\|w\| \le K$ has a solution on $[-\varepsilon, \varepsilon]$ staying in
  $U$.
\end{theorem}

\begin{proof}
  \leanok
  The Picard--Lindel\"of theorem, keeping track of the fact that the solution stays in the ball
  where the vector field is Lipschitz, gives a uniform time for $\|w\| \le \delta$. If $c$ is a
  solution, so is $t \mapsto c(st)$, with initial velocity $s\, c'(0)$; rescaling by
  $s = (K + \delta)/\delta$ gives the statement.
\end{proof}

\begin{theorem}[Uniqueness]
  \label{thm:quadratic-ode-uniqueness}
  \lean{LocallyConvex.Analysis.IsSprayCurveOn.eqOn_Ioo}
  \leanok
  \uses{thm:mb-to-frechet}
  Two solutions on $(a, b)$, staying in $U$, with the same position and velocity at some
  $t_0 \in (a, b)$, agree on $(a, b)$.
\end{theorem}

\begin{proof}
  \leanok
  Integral curves of a locally Lipschitz autonomous field agree near every point where they meet
  (Gr\"onwall); the set where they agree near $t$ is open and closed in $(a, b)$.
\end{proof}

\section{Local flows on Banach spaces}

Let $E$ be a Banach space, $U \subseteq E$ open and $f \colon U \to E$ Fr\'echet $\Ck{n}$,
$1 \le n \le \infty$. We prove that the solutions of $c' = f \circ c$ depend in a $\Ck{n}$ way on
time and initial value, following the implicit function approach of \cite[Ch.~IV, \S1]{Lan99}.

\begin{lemma}[Superposition operators]
  \label{lem:superposition-banach}
  \lean{LocallyConvex.Analysis.contDiffOn_superpose, LocallyConvex.Analysis.hasFDerivAt_superpose}
  \leanok
  Let $K$ be compact. If $f$ is $\Ck{n}$ on $U$, then $\gamma \mapsto f \circ \gamma$ is
  Fr\'echet $\Ck{n}$ on the open set of curves $\gamma \in C(K, E)$ with values in $U$ (sup norm),
  with derivative $\eta \mapsto (s \mapsto f'(\gamma(s))\,\eta(s))$.
\end{lemma}

\begin{proof}
  \leanok
  The derivative comes from the mean value inequality and the uniform continuity of $f'$ near the
  compact set $\gamma(K)$ (compare \cite[Lemma~IV.1.12]{Lan99}). It is again a superposition
  operator, composed with a continuous bilinear map, so induction on $n$ applies.
\end{proof}

\begin{theorem}[Local flows with $\Ck{n}$ dependence]
  \label{thm:banach-local-flow}
  \lean{LocallyConvex.Analysis.exists_contDiffOn_flow, LocallyConvex.Analysis.exists_smoothMBOn_flow}
  \leanok
  \uses{lem:superposition-banach}
  For $x_0 \in U$ there are $\varepsilon > 0$, an open $V \ni x_0$ and
  $\mathrm{Fl} \colon (-\varepsilon, \varepsilon) \times V \to U$, jointly $\Ck{n}$, with
  $\mathrm{Fl}(0, x) = x$ and $\partial_t \mathrm{Fl}(t, x) = f(\mathrm{Fl}(t, x))$. If $f$ is
  smooth in the sense of Michal--Bastiani, so is $\mathrm{Fl}$.
\end{theorem}

\begin{proof}
  \leanok
  Let $J\gamma(s) = \int_0^s \gamma$ on $C([-1, 1], E)$ and
  $\Psi(\lambda, x, \gamma) = (\lambda, x, \gamma - x - \lambda J(f \circ \gamma))$. By
  Lemma~\ref{lem:superposition-banach}, $\Psi$ is $\Ck{n}$, and its derivative at
  $(0, x_0, x_0)$ is a shear, hence invertible. The inverse function theorem gives a $\Ck{n}$
  family $\Gamma(\lambda, x)$ with $\Gamma = x + \lambda J(f \circ \Gamma)$; put
  $\mathrm{Fl}(t, x) = \Gamma(t, x)(1)$. By injectivity of $\Psi$ near $(0, x_0, x_0)$,
  $\Gamma(\sigma\lambda, x)(1) = \Gamma(\lambda, x)(\sigma)$ for $|\sigma| \le 1$, and the
  fundamental theorem of calculus gives the equation. Lang applies the implicit mapping theorem
  only in $x$, on a fixed short time interval \cite[Theorem~IV.1.14]{Lan99}; taking the time
  scale $\lambda$ as a parameter gives joint smoothness directly. We do not claim this argument is
  new.
\end{proof}

\begin{theorem}[Complex time]
  \label{thm:banach-local-flow-rclike}
  \lean{LocallyConvex.Analysis.exists_contDiffOn_flow_rclike}
  \leanok
  \uses{lem:superposition-banach}
  Let $\mathbb{K} \in \{\mathbb{R}, \mathbb{C}\}$, $E$ a Banach space over $\mathbb{K}$ and $f$
  of class $\Ck{n}$ over $\mathbb{K}$ on $U$, $1 \le n \le \infty$. For $x_0 \in U$ there are
  $\varepsilon > 0$, an open $V \ni x_0$ and $\mathrm{Fl} \colon B_{\mathbb{K}}(0, \varepsilon)
  \times V \to U$, jointly $\Ck{n}$ over $\mathbb{K}$, with $\mathrm{Fl}(0, x) = x$, whose
  restriction to real times solves $\partial_t \mathrm{Fl}(t, x) = f(\mathrm{Fl}(t, x))$.
\end{theorem}

\begin{proof}
  \leanok
  The proof of Theorem~\ref{thm:banach-local-flow} works verbatim over $\mathbb{K}$: the curves
  are parametrized by the real interval $[-1, 1]$ and only the time scale $\lambda$ lies in
  $\mathbb{K}$; Mathlib's inverse function theorem holds over $\mathbb{R}$ and $\mathbb{C}$. The
  rescaling identity is used for real $\sigma$ and real times. Theorem~\ref{thm:banach-local-flow}
  is the case $\mathbb{K} = \mathbb{R}$.
\end{proof}

\begin{theorem}[Holomorphic and real analytic flows {\cite[Theorem~2.4.34]{GN26}}]
  \label{thm:banach-local-flow-analytic}
  \lean{LocallyConvex.Analysis.exists_differentiableOn_flow_complex,
    LocallyConvex.Analysis.exists_realAnalyticOn_flow}
  \leanok
  \uses{thm:banach-local-flow-rclike, thm:banach-holomorphic, def:real-analytic,
    lem:normed-complexification, thm:banach-ode-uniqueness}
  If $f$ is complex differentiable on an open subset $U$ of a complex Banach space, its local
  flows are complex differentiable in complex time and initial value. If $f$ is real analytic on
  an open subset $U$ of a real Banach space, its local flows
  $\mathrm{Fl} \colon (-\varepsilon, \varepsilon) \times V \to U$ are real analytic.
\end{theorem}

\begin{proof}
  \leanok
  The holomorphic case is Theorem~\ref{thm:banach-local-flow-rclike} for $\mathbb{K} = \mathbb{C}$,
  since complex differentiable maps are $\Ck{1}$ over $\mathbb{C}$
  (Theorem~\ref{thm:banach-holomorphic}). In the real analytic case, let $g$ be a complex
  analytic extension of $f$ to an open subset of $E_{\mathbb{C}}$, a complex Banach space with
  the norm of Lemma~\ref{lem:normed-complexification}, and $\Phi$ its holomorphic flow. For real
  $t$ and $x \in E$, $t \mapsto \mathrm{Fl}(t, x)$ and $t \mapsto \Phi(t, x)$ solve
  $c' = g \circ c$ with $c(0) = x$, since $g$ extends $f$; by uniqueness they agree. So $\Phi$,
  composed with $(\mathbb{R} \times E)_{\mathbb{C}} \cong \mathbb{C} \times E_{\mathbb{C}}$, is
  a complex analytic extension of $\mathrm{Fl}$. This is the proof of Glöckner and Neeb (real
  analytic case, p.~206), except that they show that complex solutions with real data stay real
  (their Lemma~2.4.30), while we use the real flow and uniqueness.
\end{proof}

\begin{corollary}[Time-dependent equations with parameters {\cite[Theorem~2.4.34(b)]{GN26}}]
  \label{cor:banach-flow-analytic-param}
  \lean{LocallyConvex.Analysis.exists_realAnalyticOn_flow_param}
  \leanok
  \uses{thm:banach-local-flow-analytic}
  Let $E$, $\Lambda$ be real Banach spaces and $f$ real analytic on an open subset $U$ of
  $\mathbb{R} \times E \times \Lambda$. Near every $p_0 \in U$ the solution of
  $c'(t) = f(t, c(t), \lambda)$ with $c(t_0) = x$ is real analytic in $(t, t_0, x, \lambda)$.
\end{corollary}

\begin{proof}
  \leanok
  Apply Theorem~\ref{thm:banach-local-flow-analytic} to the real analytic field
  $(s, x, \lambda) \mapsto (1, f(s, x, \lambda), 0)$. Along its integral curves the first
  component has derivative $1$ and the last one derivative $0$, so the curve through
  $(t_0, x, \lambda)$ is $t \mapsto (t_0 + t, c(t_0 + t), \lambda)$. Glöckner and Neeb allow any
  locally convex parameter space; that case is Theorem~\ref{thm:lc-param-flow}.
\end{proof}

\begin{theorem}[Uniqueness]
  \label{thm:banach-ode-uniqueness}
  \lean{LocallyConvex.Analysis.eqOn_Ioo_of_contDiffOn}
  \leanok
  If $f$ is $\Ck{1}$ on $U$, two solutions of $c' = f \circ c$ on $(a, b)$ with values in $U$
  that agree at one instant agree on $(a, b)$.
\end{theorem}

\begin{proof}
  \leanok
  $f$ is locally Lipschitz; apply the argument of Theorem~\ref{thm:quadratic-ode-uniqueness}.
\end{proof}

Outside Banach spaces local existence and uniqueness can fail, even for continuous linear vector
fields \cite[Examples~2.4.2 and 2.4.3]{GN26}.

\begin{lemma}[Uniform contractions with parameters {\cite[Lemma~2.3.26(a), Proposition~2.3.27(a)]{GN26}}]
  \label{lem:uniform-contraction}
  \lean{LocallyConvex.Analysis.UniformContraction,
    LocallyConvex.Analysis.UniformContraction.continuousOn_fixedPoint,
    LocallyConvex.Analysis.UniformContraction.hasDirDerivAt_fixedPoint,
    LocallyConvex.Analysis.UniformContraction.contDiffMBOn_fixedPoint,
    LocallyConvex.Analysis.UniformContraction.dMB_fixedPoint_eq,
    LocallyConvex.Analysis.UniformContraction.contDiffMBOn_fixedPoint_of_contDiffMBOn,
    LocallyConvex.Analysis.UniformContraction.eq_fixedPoint}
  \leanok
  Let $f(p, x)$ be a family of $\theta$-contractions ($\theta < 1$) of a closed set of a Banach
  space $F$, with $p$ in an open subset of a topological vector space. The fixed point
  $\varphi(p)$ is continuous, and if $f$ is $\Ck{1}$ so is $\varphi$, with
  $\dd\varphi(p, q) = \sum_n A_p^n b_p(q)$, $A_p = \dd_2 f(p, \varphi(p))$,
  $b_p = \dd_1 f(p, \varphi(p))$. If $f$ is $\Ck{k}$, so is $\varphi$.
\end{lemma}

\begin{proof}
  \leanok
  Continuity from the uniform contraction estimate; the derivative is the Neumann series, whose
  terms converge locally uniformly (Tannery's theorem). For $\Ck{k}$, by induction on $k$:
  differentiating $\varphi(p) = f(p, \varphi(p))$, $\dd\varphi(p, q)$ is the fixed point of
  the uniform family of $\theta$-contractions $w \mapsto \dd f((p, \varphi(p)), (q, w))$ with
  parameter $(p, q)$, which is $\Ck{k-1}$ by the induction hypothesis applied to $\varphi$.
\end{proof}

\begin{lemma}[$\Ck{k}$ superposition operators with parameters {\cite[Proposition~1.7.5]{GN26}}]
  \label{lem:ck-superposition-param}
  \lean{LocallyConvex.Analysis.pairCLM,
    LocallyConvex.Analysis.contDiffMBOn_superposeParam_of_contDiffMBOn}
  \leanok
  Let $K$ be a compact metric space, $\kappa \colon K \to \mathbb{R}$ continuous, and $H$ a
  $\Ck{k}$ map on an open subset of $\mathbb{R} \times (E' \times Z)$ containing
  $\kappa(K) \times \Omega$, $E'$ normed. The operator
  $(\psi, z) \mapsto (s \mapsto H(\kappa(s), \psi(s), z))$ is $\Ck{k}$ on the open set of
  $(\psi, z) \in C(K, E') \times Z$ with $(\psi(s), z) \in \Omega$ for all $s$.
\end{lemma}

\begin{proof}
  \leanok
  The difference quotients converge uniformly in $s$ by the mean value inequality. By induction on
  $k$: the derivative is the superposition operator of $\dd H(\cdot, (0, \cdot))$, read on
  $C(K, E' \times E') \times (Z \times Z)$.
\end{proof}

\begin{theorem}[Solutions with locally convex parameters {\cite[Theorem~2.4.34]{GN26}}]
  \label{thm:lc-param-flow}
  \lean{LocallyConvex.Analysis.contDiffMBOn_superposeParam,
    LocallyConvex.Analysis.complexAnalyticOn_of_contDiffMBOn_one,
    LocallyConvex.Analysis.ODEFlowParam.complexAnalyticOn_flowPhi,
    LocallyConvex.Analysis.ODEFlowParam.hasDerivAt_flowPhi_time,
    LocallyConvex.Analysis.ODEFlowParam.flowPhi_mem_of_real,
    LocallyConvex.Analysis.exists_complexAnalyticOn_flow_param,
    LocallyConvex.Analysis.exists_realAnalyticOn_flow_lcParam,
    LocallyConvex.Analysis.ODEFlowParam.contDiffMBOn_flowPhi_of_contDiffMBOn,
    LocallyConvex.Analysis.ODEFlowParam.hasDerivAt_flowPhi_time_of,
    LocallyConvex.Analysis.exists_contDiffMBOn_flow_lcParam,
    LocallyConvex.Analysis.ODEFlowParam.exists_bounds_of_lipschitz,
    LocallyConvex.Analysis.ODEFlowParam.hasDerivAt_flowPhi_time_continuous,
    LocallyConvex.Analysis.exists_continuousOn_flow_lcParam}
  \leanok
  \uses{lem:uniform-contraction, lem:ck-superposition-param, def:real-analytic,
    lem:normed-complexification}
  Let $E$ be a complex Banach space, $W$ a complex locally convex space and $f$ complex analytic on
  an open subset of $\mathbb{C} \times E \times W$. The solution $\Phi(t, t_0, y_0, p)$ of
  $y' = f(t, y, p)$, $y(t_0) = y_0$, is complex analytic in all variables near any point. The same
  holds for real analytic $f$ on $\mathbb{R} \times E \times W$, $E$ a real Banach space and $W$ a
  real locally convex space, and if $f$ is $\Ck{k}$, $k \ge 1$, the solution is $\Ck{k}$ in all
  variables; if $f$ is continuous and Lipschitz in $y$, uniformly in the other variables near
  the point, the solution is continuous.
\end{theorem}

\begin{proof}
  \leanok
  As in \cite{GN26}: $\eta(\tau) = \Phi(t_0 + \tau(t - t_0), \ldots) - y_0$ is the fixed point
  of $\eta \mapsto \int_0^{\tau} (t - t_0) f(t_0 + s(t - t_0), \eta(s) + y_0, p)\,ds$ on
  $C([-1, 1], E)$, a uniform family of contractions, $\Ck{1}$ by the superposition results; its
  differential is complex linear, and a real $\Ck{1}$ map with complex linear differentials into a
  Banach space is complex analytic \cite[Theorem~2.1.12]{GN26}. The rescaling identity (2.107)
  gives the differential equation. In the real case, the complexified equation preserves the closed
  set of curves with values in the real subspace, so its fixed point is real for real data; this
  replaces \cite[Lemma~2.4.30]{GN26}. In the $\Ck{k}$ case the map $g$ is $\Ck{k}$ by
  Lemma~\ref{lem:ck-superposition-param}, hence so is its fixed point
  (Lemma~\ref{lem:uniform-contraction}); the construction is written once for $\mathbb{K} \in
  \{\mathbb{R}, \mathbb{C}\}$. For continuous $f$, the bounds come from continuity and the
  Lipschitz condition, and $\Phi(r) = y_0 + \int_{t_0}^r f(v, \Phi(v), p)\,dv$.
\end{proof}

\begin{theorem}[Maximal flows {\cite[Theorem~2.4.29, Lemma~2.5.7, Theorem~2.5.18, Remark~2.4.35]{GN26}}]
  \label{thm:maximal-flow}
  \lean{LocallyConvex.Analysis.MaximalFlow.dom, LocallyConvex.Analysis.MaximalFlow.flow,
    LocallyConvex.Analysis.MaximalFlow.IsSol.eqOn, LocallyConvex.Analysis.MaximalFlow.flow_eq,
    LocallyConvex.Analysis.MaximalFlow.flow_concat, LocallyConvex.Analysis.MaximalFlow.isSol_flow,
    LocallyConvex.Analysis.MaximalFlow.FlowRegularity,
    LocallyConvex.Analysis.MaximalFlow.HasLocalFlows,
    LocallyConvex.Analysis.MaximalFlow.isOpen_dom_and,
    LocallyConvex.Analysis.MaximalFlow.locallyLipschitzY_of_realAnalyticOn,
    LocallyConvex.Analysis.MaximalFlow.isOpen_dom_and_realAnalyticOn_flow,
    LocallyConvex.Analysis.MaximalFlow.flowRegularity_contDiffMBOn,
    LocallyConvex.Analysis.MaximalFlow.locallyLipschitzY_of_contDiffMBOn,
    LocallyConvex.Analysis.MaximalFlow.isOpen_dom_and_contDiffMBOn_flow,
    LocallyConvex.Analysis.MaximalFlow.isOpen_dom_and_smoothMBOn_flow,
    LocallyConvex.Analysis.MaximalFlow.LocallyLipschitzYParam,
    LocallyConvex.Analysis.MaximalFlow.flowRegularity_continuousOn,
    LocallyConvex.Analysis.MaximalFlow.isOpen_dom_and_continuousOn_flow,
    LocallyConvex.Analysis.MaximalFlow.realAnalyticOn_flow_of_forall_isSol}
  \leanok
  \uses{thm:lc-param-flow}
  Let $E$ be a real Banach space, $W$ a topological space and $f$ defined on
  $U \subseteq \mathbb{R} \times E \times W$, locally Lipschitz in $y$. Solutions of
  $y' = f(t, y, p)$ on open intervals are unique and glue to maximal solutions; let
  $\mathrm{Fl}(t, t_0, y_0, p)$ be the maximal flow, on its domain $\Omega$. Then
  $\mathrm{Fl}(t, t_2, \mathrm{Fl}(t_2, t_1, y_1, q), q) = \mathrm{Fl}(t, t_1, y_1, q)$. If the
  equation admits local flows of a regularity class (closed under restriction, gluing and the
  composition below), $\Omega$ is open and $\mathrm{Fl}$ has that regularity on it. In particular,
  if $W$ is a real locally convex space and $f$ is continuous and locally Lipschitz in $y$ (with
  parameters), $\Ck{k}$ ($k \ge 1$), smooth or real analytic on an open set $U$, then $\Omega$ is
  open and $\mathrm{Fl}$ has the same regularity on $\Omega$; if all solutions with initial data in
  $J \times S$ exist on the open interval $J$, $\mathrm{Fl}$ is real analytic on $J \times J \times S$.
\end{theorem}

\begin{proof}
  \leanok
  Uniqueness on open intervals by connectedness, from local uniqueness (Gronwall). Globalization as
  in \cite[Theorem~2.5.18]{GN26}: given $(\tau, t_0, y_0, p) \in \Omega$, the set of times $t_1$
  near which $\Omega$ is a neighbourhood and $\mathrm{Fl}$ is regular contains $t_0$ (a local flow)
  and is open and closed in the interval of the solution: near a limit point $\theta$, with a
  regular $t_2$ close to $\theta$, $\mathrm{Fl}(t, t_1, y_1, q) =
  \Phi_\theta(t, t_2, \mathrm{Fl}(t_2, t_1, y_1, q), q)$ for the local flow $\Phi_\theta$ at
  $\theta$. A real analytic $f$ is $\Ck{1}$ in $(t, y)$ for fixed $p$, hence locally Lipschitz in
  $y$, and admits real analytic local flows by Theorem~\ref{thm:lc-param-flow}; a $\Ck{k}$ map,
  $k \ge 1$, is locally Lipschitz in $y$ by the continuity of $\dd f$, and admits $\Ck{k}$ local
  flows by the same theorem; so does a continuous map with a local Lipschitz condition.
\end{proof}



\section{The local Frobenius theorem on Banach spaces}

Let $\mathbb{K} \in \{\mathbb{R}, \mathbb{C}\}$, $E$, $F$ Banach spaces over $\mathbb{K}$,
$O \subseteq E \times F$ open and $f \colon O \to L(E, F)$ of class $\Ck{n}$ over $\mathbb{K}$,
$1 \le n \le \infty$. We follow Penot, who derives the Frobenius theorem from the implicit
function theorem \cite[Théorème~1.1]{Pen70}.

\begin{lemma}[Penot's map]
  \label{lem:penot-map}
  \lean{LocallyConvex.Analysis.contDiffOn_penotMap, LocallyConvex.Analysis.hasFDerivAt_penotRHS,
    LocallyConvex.Analysis.hasFDerivAt_penotMap_zero}
  \leanok
  \uses{lem:superposition-banach}
  The map $P(x, y, \delta) = \bigl(x, y, \delta - (s \mapsto f(sx, y + \int_0^s \delta)\,x)\bigr)$
  on $E \times F \times C([-1, 1], F)$ is $\Ck{n}$ where the curve
  $s \mapsto (sx, y + \int_0^s \delta)$ has values in $O$. If $(0, y_0) \in O$, its derivative at
  $(0, y_0, 0)$ is the invertible shear $(a, b, c) \mapsto (a, b, c - f(0, y_0)\,a)$.
\end{lemma}

\begin{proof}
  \leanok
  $P$ is built from continuous linear maps, the superposition operator of $f$
  (Lemma~\ref{lem:superposition-banach}) and the continuous bilinear pointwise application
  $C([-1,1], L(E, F)) \times C([-1,1], E) \to C([-1,1], F)$. At $x = 0$ the term
  $f(\dots)\,x$ vanishes, and so does the part of its derivative that differentiates $f$.
\end{proof}

\begin{theorem}[Local Frobenius theorem {\cite[Théorème~1.1 and Remarque~1]{Pen70}}]
  \label{thm:frobenius-banach}
  \lean{LocallyConvex.Analysis.exists_frobenius, LocallyConvex.Analysis.exists_frobenius_zero,
    LocallyConvex.Analysis.exists_frobenius_holomorphic}
  \leanok
  \uses{lem:penot-map, thm:banach-holomorphic}
  Assume (Frob): for $q \in O$ and $a, b \in E$,
  \[ D_1 f(q)(a)(b) + D_2 f(q)(f(q)\,a)(b) = D_1 f(q)(b)(a) + D_2 f(q)(f(q)\,b)(a). \]
  For every $(x_0, y_0) \in O$ there are $r > 0$ with $B(x_0, r) \times B(y_0, r) \subseteq O$ and
  a $\Ck{n}$ map $\alpha \colon B(x_0, r) \times B(y_0, r) \to F$ with $\alpha(x_0, y) = y$ and
  $D_1 \alpha(x, y) = f(x, \alpha(x, y))$. If $\mathbb{K} = \mathbb{C}$ and $f$ is holomorphic,
  $\alpha$ is holomorphic.
\end{theorem}

\begin{proof}
  \leanok
  Translate $x_0$ to $0$. By Lemma~\ref{lem:penot-map} and the inverse function theorem there is a
  $\Ck{n}$ family $\varphi(x, y) \in C([-1, 1], F)$ with $P(x, y, \varphi(x, y)) = (x, y, 0)$ and
  $\varphi(0, y) = 0$; put $\alpha(x, y) = y + \int_0^1 \varphi(x, y)$. Let
  $c(s) = (sx, y + \int_0^s \varphi(x, y))$ and $z \in E$. Differentiating the fixed point
  identity, $K(s) = \partial_x\bigl(\int_0^s \varphi\bigr)(x, y)\,z$ solves the linear equation
  $u' = Df(c(s))(sz, u)\,x + f(c(s))\,z$ with $u(0) = 0$. By (Frob), so does
  $L(s) = s\, f(c(s))\,z$, and uniqueness for linear equations gives $K(1) = L(1)$, which is
  $D_1 \alpha(x, y)\,z = f(x, \alpha(x, y))\,z$. Penot instead shows that $D_3 P$ maps $K'$ and
  $L'$ to the same point and that it is invertible near $0$; and he works with curves of class
  $\Ck{1}$ and the implicit function theorem, where we take $\delta = \gamma'$ as unknown. We do
  not claim this argument is new. For $\mathbb{K} = \mathbb{C}$, holomorphic maps are $\Ck{1}$
  over $\mathbb{C}$ and $\Ck{1}$ maps over $\mathbb{C}$ are holomorphic
  (Theorem~\ref{thm:banach-holomorphic}).
\end{proof}

\begin{theorem}[Local Frobenius theorem, real analytic case {\cite[Théorème~1.1]{Pen70}}]
  \label{thm:frobenius-real-analytic}
  \lean{LocallyConvex.Analysis.exists_frobenius_realAnalytic,
    LocallyConvex.Analysis.frobeniusRayField,
    LocallyConvex.Analysis.realAnalyticOn_frobeniusRayField}
  \leanok
  \uses{thm:frobenius-banach, thm:banach-local-flow-analytic, thm:banach-ode-uniqueness,
    def:real-analytic}
  Let $E$, $F$ be real Banach spaces and $f \colon O \to L(E, F)$ real analytic in the sense of
  Milnor, satisfying (Frob). Then the map $\alpha$ of Theorem~\ref{thm:frobenius-banach} is real
  analytic on some $B(x_0, r) \times B(y_0, r)$.
\end{theorem}

\begin{proof}
  \leanok
  Penot uses a real analytic implicit function theorem; we argue with flows instead. The vector
  field $X(s, v, c) = (1, 0, f(x_0 + sv, c)\,v)$ on $\mathbb{R} \times E \times F$ is real
  analytic. Since $\frac{d}{dt}\alpha(x_0 + tv, y) = f(x_0 + tv, \alpha(x_0 + tv, y))\,v$, the
  curve $t \mapsto (t, v, \alpha(x_0 + tv, y))$ is an integral curve of $X$ through $(0, v, y)$,
  so by uniqueness it is $t \mapsto \mathrm{Fl}(t, (0, v, y))$, where $\mathrm{Fl}$ is the local
  flow of $X$ on $(-\varepsilon, \varepsilon) \times V$. Hence $\alpha(x, y)$ is the last
  component of $\mathrm{Fl}(\varepsilon/2, (0, \tfrac{2}{\varepsilon}(x - x_0), y))$, which is
  real analytic by Theorem~\ref{thm:banach-local-flow-analytic}. We do not claim this argument is
  new.
\end{proof}

\begin{proposition}[Involutivity in local form {\cite[Chapter~VI, \S1, p.~157]{Lan99}}]
  \label{prop:frobenius-involutive}
  \lean{LocallyConvex.Analysis.lieBracket_mem_graph_iff,
    LocallyConvex.Analysis.frob_iff_involutive}
  \leanok
  Let $E$, $F$ be normed spaces and $f \colon E \times F \to L(E, F)$ differentiable at $q$. If
  the vector fields $X$, $Y$ are differentiable at $q$ and $X(p), Y(p) \in \operatorname{graph}
  f(p)$ for $p$ near $q$, then $[X, Y](q) \in \operatorname{graph} f(q)$ if and only if
  $Df(q)(X(q))(Y_1(q)) = Df(q)(Y(q))(X_1(q))$. Hence (Frob) at $q$ holds if and only if the
  distribution $p \mapsto \operatorname{graph} f(p)$ is involutive at $q$.
\end{proposition}

\begin{proof}
  \leanok
  Differentiating $X_2 = f\,X_1$ gives $DX_2(q)\,v = Df(q)(v)\,X_1(q) + f(q)\,DX_1(q)\,v$, and
  similarly for $Y$. In $[X, Y](q)_2 - f(q)\,[X, Y](q)_1$ the terms with $DX_1$ and $DY_1$ cancel,
  leaving $Df(q)(X(q))\,Y_1(q) - Df(q)(Y(q))\,X_1(q)$. For the converse take the fields
  $p \mapsto (a, f(p)\,a)$ and $p \mapsto (b, f(p)\,b)$.
\end{proof}

\begin{theorem}[Frobenius theorem in chart form {\cite[Chapter~VI, Theorems~1.1 and 1.2]{Lan99}}]
  \label{thm:frobenius-chart}
  \lean{LocallyConvex.Analysis.exists_frobenius_chart,
    LocallyConvex.Analysis.exists_frobenius_chart_holomorphic,
    LocallyConvex.Analysis.map_range_inl_eq_graph}
  \leanok
  \uses{thm:frobenius-banach, prop:frobenius-involutive}
  Under the hypotheses of Theorem~\ref{thm:frobenius-banach}, every $(x_0, y_0) \in O$ is fixed
  by a $\Ck{n}$ diffeomorphism $\Phi(x, y) = (x, \alpha(x, y))$ between open subsets of $O$ with
  $D\Phi(p)(a, 0) = (a, f(\Phi(p))\,a)$. Thus $D\Phi^{-1}(z)$ maps $\operatorname{graph} f(z)$
  onto $E \times 0$: the chart $\Phi^{-1}$ rectifies the distribution. If $\mathbb{K} = \mathbb{C}$
  and $f$ is holomorphic, $\Phi$ and $\Phi^{-1}$ are holomorphic.
\end{theorem}

\begin{proof}
  \leanok
  Take $\alpha$ from Theorem~\ref{thm:frobenius-banach}. Since $D_1\alpha = f(x, \alpha)$ and
  $\alpha(x_0, \cdot) = \mathrm{id}$, the derivative of $\Phi$ at $(x_0, y_0)$ is the shear
  $(a, b) \mapsto (a, f(x_0, y_0)\,a + b)$, which is invertible. Apply the inverse function
  theorem and restrict to the open set where $D\Phi$ is invertible.
\end{proof}

\section{The Nijenhuis tensor}

\begin{definition}[Nijenhuis tensor {\cite[p.~73]{Pen70}}]
  \label{def:nijenhuis}
  \lean{LocallyConvex.Analysis.nijenhuis, LocallyConvex.Analysis.nijenhuis_apply_eq}
  \leanok
  For a field of operators $A \colon G \to L(G)$ differentiable at $z$, put
  \[ N_A(z)(a, b) = DA(z)(A a)\,b - DA(z)(A b)\,a + A\,DA(z)(b)\,a - A\,DA(z)(a)\,b, \]
  with $A = A(z)$. For vector fields $X$, $Y$ differentiable at $z$,
  $N_A(z)(X(z), Y(z)) = \bigl([AX, AY] + A^2[X, Y] - A[X, AY] - A[AX, Y]\bigr)(z)$.
\end{definition}

\begin{theorem}[Vanishing Nijenhuis tensor {\cite[Théorème~3.5, (4) $\Rightarrow$ (6)]{Pen70}}]
  \label{thm:nijenhuis-involutive}
  \lean{LocallyConvex.Analysis.lieBracket_mem_ker_of_nijenhuis,
    LocallyConvex.Analysis.fderiv_anticomm_of_sq_eq_neg_one}
  \leanok
  \uses{def:nijenhuis}
  Let $G$ be a complex normed space and $F$ differentiable at $z$ with $F^2 = -1$ near $z$ and
  $N_F(z) = 0$. Then $\ker(F + i)$ is involutive at $z$.
\end{theorem}

\begin{proof}
  \leanok
  Differentiating $F^2 = -1$, $DF(z)(v)$ anticommutes with $F(z)$ and maps $\ker(F + i)$ into
  $\ker(F - i)$. For $Z$, $W$ of type $(0, 1)$, $c = DF(W)Z - DF(Z)W$ lies in $\ker(F - i)$, while
  $N_F(Z, W) = (F + i)c$ puts it in $\ker(F + i)$; so $c = 0$. Differentiating $FZ = -iZ$ gives
  $(F + i)[Z, W] = c$.
\end{proof}

\begin{proposition}[Complexification {\cite[Théorème~3.5, (2) $\Rightarrow$ (4)]{Pen70}}]
  \label{prop:nijenhuis-complexified}
  \lean{LocallyConvex.Analysis.nijenhuis_ofReal,
    LocallyConvex.Analysis.nijenhuis_eq_zero_of_ofReal,
    LocallyConvex.Analysis.fderiv_ofReal_of_extends}
  \leanok
  \uses{def:nijenhuis}
  If $F(x + i0)$ extends $J(x)$ complex linearly for real $x$ near $x_0$, then
  $N_F(x_0)(a + i0, b + i0) = N_J(x_0)(a, b) + i0$; hence $N_J(x_0) = 0$ implies $N_F(x_0) = 0$.
\end{proposition}

\begin{proof}
  \leanok
  The derivative of $F$ in real directions extends that of $J$, and a complex bilinear map on
  $E_{\mathbb{C}}$ vanishing on real vectors vanishes.
\end{proof}

\begin{theorem}[Graph form of the $(0, 1)$ distribution {\cite[p.~75]{Pen70}}]
  \label{thm:nijenhuis-graph}
  \lean{LocallyConvex.Analysis.exists_graph_frob_of_nijenhuis,
    LocallyConvex.Analysis.apply_eq_neg_I_iff_graph, LocallyConvex.Analysis.typeSplit,
    LocallyConvex.Analysis.lieBracket_comp_equiv}
  \leanok
  \uses{thm:nijenhuis-involutive, prop:frobenius-involutive}
  Let $G$ be a complex Banach space and $F \colon O \to L(G)$ holomorphic with $F^2 = -1$ and
  $N_F = 0$. Near $z_0 \in O$, in the coordinates $G \cong \ker(F(z_0) + i) \times
  \ker(F(z_0) - i)$, the distribution $z \mapsto \ker(F(z) + i)$ is the field of graphs of a
  holomorphic map satisfying (Frob).
\end{theorem}

\begin{proof}
  \leanok
  With $Q = (1 + iF(z_0))/2$, $P = 1 - Q$ and $Q_z = (1 + iF(z))/2$, the space $\ker(F(z) + i)$
  is the graph of $T_z S_z^{-1}$, $S_z = Q Q_z$, $T_z = P Q_z$ on $\ker(F(z_0) + i)$, as long as
  $S_z$ is invertible and $\|Q_z - Q\| < 1$. (Frob) follows from
  Theorem~\ref{thm:nijenhuis-involutive} and Proposition~\ref{prop:frobenius-involutive}. Penot
  works with the projector and his Théorème~2.3 instead; we do not claim this argument is new.
\end{proof}

\section{The Newlander--Nirenberg theorem, real analytic case}

\begin{lemma}[Holomorphic extension of a real analytic field of operators]
  \label{lem:holomorphic-extension-clm}
  \lean{LocallyConvex.Analysis.RealAnalyticOn.exists_holomorphic_extension_clm,
    LocallyConvex.Analysis.complexifyCLM, LocallyConvex.Analysis.extendC,
    LocallyConvex.Analysis.clmOfComplexification}
  \leanok
  If $J \colon E \to L(E)$ is real analytic on $U$, there is a holomorphic
  $F \colon W \to L_{\mathbb{C}}(E_{\mathbb{C}})$ on an open $W \supseteq U + i0$ with
  $F(x + i0) = J(x)_{\mathbb{C}}$.
\end{lemma}

\begin{proof}
  \leanok
  Compose the complex analytic extension of $J$ given by Milnor's definition with the continuous
  complex linear map $A + iB \mapsto (u + iv \mapsto (Au - Bv) + i(Av + Bu))$.
\end{proof}

\begin{theorem}[Real analytic inverse function theorem]
  \label{thm:real-analytic-inverse}
  \lean{LocallyConvex.Analysis.RealAnalyticOn.exists_openPartialHomeomorph,
    OpenPartialHomeomorph.realAnalyticOn_symm}
  \leanok
  Let $E$, $F$ be real Banach spaces. A real analytic map $f \colon U \to F$ (in the sense of
  Milnor) with $Df(x_0)$ invertible restricts to a real analytic diffeomorphism between open
  neighbourhoods of $x_0$ and $f(x_0)$; the inverse of a real analytic diffeomorphism is real
  analytic.
\end{theorem}

\begin{proof}
  \leanok
  A complex analytic extension $g$ of $f$ to an open subset of the normed complexification
  $E_{\mathbb{C}}$ has derivative $Df(x_0)_{\mathbb{C}}$ at $x_0 + i0$, which is invertible. The
  holomorphic inverse function theorem gives a holomorphic local inverse $\psi$ of $g$, and
  $\psi(y + i0) = f^{-1}(y) + i0$ for real $y$ near $f(x_0)$, since
  $g(f^{-1}(y) + i0) = y + i0$. So $\psi$ extends $f^{-1}$ near $f(x_0)$, and real analyticity is
  a local property \cite[Lemma~2.2.3]{GN26}.
\end{proof}

\begin{theorem}[Local Newlander--Nirenberg theorem {\cite[Théorème~3.5, (6) $\Rightarrow$ (1), pp.~75--76]{Pen70}}]
  \label{thm:newlander-nirenberg}
  \lean{LocallyConvex.Analysis.exists_holomorphic_chart,
    LocallyConvex.Analysis.exists_transverse_coordinate,
    LocallyConvex.Analysis.eqOn_ball_of_ofReal}
  \leanok
  \uses{lem:holomorphic-extension-clm, prop:nijenhuis-complexified, thm:nijenhuis-graph,
    thm:frobenius-chart, thm:real-analytic-inverse}
  Let $E$ be a real Banach space, $U \subseteq E$ open and $J \colon U \to L(E)$ real analytic
  with $J^2 = -1$ and $N_J = 0$. Every $x_0 \in U$ lies in the domain of a chart $\varphi$, real
  analytic with real analytic inverse, into the complex Banach space $\ker(J(x_0)_{\mathbb{C}} - i)$, with
  $D\varphi(x)\,J(x) = i\,D\varphi(x)$.
\end{theorem}

\begin{proof}
  \leanok
  Extend $J$ to $F$ (Lemma~\ref{lem:holomorphic-extension-clm}). $F^2 = -1$ and $N_F = 0$ hold at
  the real points (Proposition~\ref{prop:nijenhuis-complexified}), hence on a ball around
  $x_0 + i0$ by the identity theorem. By Theorem~\ref{thm:nijenhuis-graph} and the holomorphic
  Frobenius theorem (Theorem~\ref{thm:frobenius-chart}) there is $\Phi$ with $D\Phi^{-1}$ mapping
  $\ker(F + i)$ into $\ker(F(z_0) + i) \times 0$. The transverse coordinate
  $w = (\Phi^{-1} \circ \ell)_2$ kills $\ker(F + i)$; at real points, applied to
  $a + i J(x) a \in \ker(F + i)$, this gives $Dw(x)(J(x) a) = i\,Dw(x)\,a$. At $x_0$ the kernel
  of $Dw$ is $\ker(F(z_0) + i)$, which meets $E$ only in $0$, and $Dw$ is onto, so
  $D(w|_E)(x_0)$ is a real isomorphism and the inverse function theorem applies. The chart
  $w|_E$ is real analytic since $w$ is holomorphic, and so is its inverse
  (Theorem~\ref{thm:real-analytic-inverse}). Penot uses
  $x \mapsto f(x, 0)$ instead; we do not claim this variant is new.
\end{proof}

\begin{theorem}[Complex structures]
  \label{thm:complex-structure}
  \lean{LocallyConvex.Analysis.WithComplexStructure,
    LocallyConvex.Analysis.WithComplexStructure.equivOfIntertwining,
    LocallyConvex.Analysis.WithComplexStructure.nonempty_equiv_of_norm_lt,
    LocallyConvex.Analysis.WithComplexStructure.equivTypeOneZero}
  \leanok
  \uses{thm:newlander-nirenberg}
  A complex structure $J$ ($J^2 = -1$) on a real topological vector space $E$ makes it a complex
  topological vector space, $(a + ib)v = av + bJv$. A real isomorphism $S$ with $J'S = SJ$ is a
  complex isomorphism $(E, J) \cong (E, J')$; on a Banach space, if $\|(J' - J)J\| < 2$ then
  $S = 1 - J'J$ is such an isomorphism. Finally $(E, J) \cong \ker(J_{\mathbb{C}} - i)$ by
  $v \mapsto (v - iJv)/2$: the models of the holomorphic charts of
  Theorem~\ref{thm:newlander-nirenberg} at different points are isomorphic for nearby points.
\end{theorem}

\begin{proof}
  \leanok
  $J'S = J' + J = SJ$, and $S = 2 - (J' - J)J$ is invertible by the Neumann series. The inverse of
  $v \mapsto (v - iJv)/2$ is $w \mapsto 2\,\mathrm{Re}\,w$.
\end{proof}

